
\section{Disegno degli esperimenti}

La campagna \`e articolata in tre esperimenti, ciascuno costruito per isolare un solo fattore
tenendo fissi tutti gli altri. Complessivamente sono state eseguite 485 esecuzioni misurate,
riportate in Tabella~\ref{tab:doe}.

\begin{table}[h]\centering\small
\begin{tabular}{llccl}
\toprule
& \textbf{fattore studiato} & \textbf{celle} & \textbf{esecuzioni} & \textbf{carico} \\
\midrule
Esperimento 1 & granularit\`a del tracciamento & 4 & 80 & solo task critico \\
Esperimento 2 & criterio di campionamento & 6 & 150 & 1 critico + 4 di disturbo \\
Esperimento 3 & modalit\`a di consegna $\times$ carico & 12 & 180 & 0, 1, 4, 8 di disturbo \\
\addlinespace
verifica & controllo di riproducibilit\`a & 3 & 75 & 1 critico + 4 di disturbo \\
\bottomrule
\end{tabular}
\caption{Struttura della campagna sperimentale.}
\label{tab:doe}
\end{table}

\subsection{Esperimento 1 --- costo puro della strumentazione}

\textbf{Obiettivo.} Quantificare quanto costa al task critico la sola presenza della
strumentazione, in assenza di qualunque altro carico.

\textbf{Costruzione.} Un solo task periodico critico, nessun task di disturbo, campionamento sempre
attivo e consegna differita. L'unico fattore che varia \`e la granularit\`a del tracciamento sui
suoi quattro livelli, con 20 ripetizioni per livello. L'assenza di carico concorrente \`e voluta:
serve a ottenere un valore di riferimento non contaminato da contesa, rispetto al quale
interpretare gli esperimenti successivi.

\subsection{Esperimento 2 --- il campionamento protegge il task critico?}

\textbf{Obiettivo.} \`E la domanda centrale dello studio. In un sistema a criticit\`a mista ci si
attende di poter conservare integralmente le tracce dei task critici riducendo quelle dei task
best-effort. L'esperimento verifica se il criterio di campionamento proporzionale offerto dalla
libreria consenta questa distinzione.

\textbf{Costruzione.} Carico misto realistico (un task critico e quattro di disturbo), granularit\`a
e modalit\`a di consegna fissate, e sei livelli del solo criterio di campionamento: sempre attivo,
mai attivo, e quattro proporzioni intermedie (0.1, 0.3, 0.5, 0.7), con 25 ripetizioni ciascuno. Le
25 ripetizioni sono necessarie perch\'e l'esito \`e probabilistico e la grandezza di interesse \`e
una frequenza, di cui occorre stimare l'incertezza.

\textbf{Variabile osservata.} A differenza degli altri due esperimenti, qui non interessa
principalmente il tempo, ma \emph{quanti} elementi di traccia vengono effettivamente emessi e a
quale task appartengono. Per questo l'esperimento adotta la modalit\`a di uscita che li rende
contabili.

\subsection{Esperimento 3 --- consegna e contesa sotto carico crescente}

\textbf{Obiettivo.} Verificare come il costo della telemetria si comporti al crescere del numero di
task concorrenti, e soprattutto se la modalit\`a con cui gli elementi vengono consegnati incida
sulla capacit\`a del task critico di rispettare le proprie scadenze.

\textbf{Costruzione.} Due fattori incrociati: la modalit\`a di consegna (differita o immediata) e il
numero di task di disturbo (0, 1, 4, 8), con 15 ripetizioni per combinazione. A questi si affianca,
per ogni livello di carico, una condizione di controllo con tracciamento disattivato, che fornisce
il riferimento rispetto al quale misurare il costo a parit\`a di contesa. \`E l'esperimento pi\`u
vicino a una situazione d'esercizio.

\subsection{Scelte metodologiche comuni}

\textbf{Numero di ripetizioni.} Da 15 a 25 per cella, in linea con quanto richiesto per garantire
significativit\`a statistica, e dimensionato sulla variabilit\`a osservata nelle prove preliminari.

\textbf{Ordine di esecuzione interlacciato.} Le ripetizioni non sono eseguite raggruppate per cella,
ma alternando le celle a ogni ripetizione. Eseguire tutte le ripetizioni di una configurazione e
poi quelle della successiva avrebbe fatto s\`i che qualunque deriva lungo la campagna --- termica o
di altra natura --- si sommasse a una sola configurazione, diventando indistinguibile dal fattore in
studio. Con l'ordine interlacciato, un'eventuale deriva agisce su tutte le celle allo stesso modo.

\textbf{Verifiche preliminari bloccanti.} Prima di ogni campagna si controlla automaticamente che
l'isolamento delle CPU sia attivo e che la frequenza sia effettivamente fissata; in caso contrario
l'esecuzione non parte. La cautela nasce da un'esperienza concreta: in assenza di isolamento la
procedura di lancio ricade silenziosamente su un'esecuzione non isolata, e un'intera campagna
risulterebbe inutilizzabile senza alcuna segnalazione.

\textbf{Isolamento dei dati per esecuzione.} Ogni ripetizione scrive in una propria cartella, che
contiene anche la configurazione effettivamente eseguita. Questo evita la sovrascrittura dei log
fra esecuzioni successive e rende ciascun risultato verificabile a posteriori.

\textbf{Registrazione delle condizioni ambientali.} A ogni esecuzione vengono registrate la
frequenza osservata e la temperatura prima e dopo, in modo che un'eventuale deriva della
piattaforma sia rilevabile nei dati anzich\'e restare implicita.

\textbf{Trattamento del transitorio iniziale.} Le prime iterazioni di ogni esecuzione presentano un
margine negativo per costruzione: il task critico fissa l'istante di riferimento e attende poi che
tutti gli altri thread siano pronti, cosicch\'e alla ripartenza l'origine dei tempi risulta gi\`a
arretrata. Il numero di iterazioni interessate \textbf{cresce con il numero di thread} (da 1 con il
solo task critico a circa 10 con otto task di disturbo). Poich\'e il numero di task di disturbo \`e
uno dei fattori dell'Esperimento 3, scartare un numero fisso di righe avrebbe introdotto un errore
sistematico correlato proprio con il fattore in studio, che si sarebbe letto come ``pi\`u carico,
pi\`u scadenze mancate''. L'analisi scarta quindi le righe iniziali fino al primo margine non
negativo, registrando quante ne sono state escluse; oltre quel punto, un margine negativo \`e una
scadenza realmente mancata.

\section{Esecuzione degli esperimenti}

Questa sezione riporta la sequenza di comandi che porta dalla macchina appena avviata ai dati
analizzabili. Gli script citati sono quelli sviluppati per la campagna.

\subsection{Preparazione della piattaforma}

Da ripetere a ogni sessione, poich\'e nessuna delle due impostazioni sopravvive al riavvio.

\begin{lstlisting}
# 1. frequenza fissata allo stato di prestazione massimo
sudo ./scripts/utils_isolation/pin_cpu_freq.sh fix 0
sudo ./scripts/utils_isolation/pin_cpu_freq.sh status   # verifica

# 2. isolamento dei due core fisici destinati agli esperimenti
sudo ./scripts/utils_isolation/isolate_cpus.sh 2,3,6,7

# al termine della sessione
sudo ./scripts/utils_isolation/reset_isolation.sh
\end{lstlisting}

\subsection{Definizione del workload}

La configurazione viene generata a partire dai parametri del carico, cos\`i che le celle differiscano
soltanto per ci\`o che si intende variare:

\begin{lstlisting}
python3 scripts/measurements/gen_config.py \
        --n-lo 4              `# numero di task di disturbo` \
        --duration 20         `# durata in secondi` \
        --hi-cpu 2 --lo-cpus 6  `# core fisici distinti` \
        --calibration 29 \
        --out 1-configs/cfg_n4.json
\end{lstlisting}

Lo script avverte se il task critico e quelli di disturbo dovessero ricadere sui due thread dello
stesso core fisico, condizione da evitare per le ragioni esposte in \S1.1.

\subsection{Definizione del tipo di monitoraggio}

Poich\'e i parametri della telemetria sono fissati in compilazione, ogni configurazione corrisponde
a un eseguibile distinto:

\begin{lstlisting}
# esempio: granularita' massima, consegna immediata, campionamento sempre attivo
make -C rt-app/src CPPFLAGS="-DRTAPP_TRACE_LEVEL=3 \
                             -DRTAPP_PROCESSOR_TYPE=1 \
                             -DRTAPP_SAMPLER_TYPE=0 \
                             -DRTAPP_SAMPLER_RATIO=0.0 \
                             -DRTAPP_EXPORTER_TYPE=0"

# esempio: campionamento proporzionale al 30%, elementi contabili sull'uscita standard
make -C rt-app/src CPPFLAGS="-DRTAPP_TRACE_LEVEL=2 \
                             -DRTAPP_PROCESSOR_TYPE=0 \
                             -DRTAPP_SAMPLER_TYPE=1 \
                             -DRTAPP_SAMPLER_RATIO=0.3 \
                             -DRTAPP_EXPORTER_TYPE=1"
\end{lstlisting}

Gli eseguibili prodotti vengono conservati in una cache indicizzata dalla combinazione di
parametri, in modo che ogni cella dell'esperimento usi esattamente il binario corrispondente e che
la compilazione non venga ripetuta inutilmente.

\subsection{Esecuzione di una singola prova e di una campagna}

\begin{lstlisting}
# singola esecuzione, all'interno dell'area isolata
bash scripts/measurements/test.sh <cartella_di_uscita> <eseguibile> <configurazione>

# campagna completa: costruisce gli eseguibili, genera le configurazioni,
# esegue le ripetizioni in ordine interlacciato e registra le condizioni ambientali
./scripts/measurements/run_doe.sh block1     # granularita' del tracciamento
./scripts/measurements/run_doe.sh block2     # criterio di campionamento
./scripts/measurements/run_doe.sh block3     # modalita' di consegna e carico
\end{lstlisting}

Ogni esecuzione produce la configurazione effettivamente usata, un log per ciascun thread, e le
uscite standard e di errore. Al termine, una riga per esecuzione viene aggiunta a una tabella
riassuntiva con i livelli dei fattori e le condizioni ambientali registrate.

\subsection{Analisi}

\begin{lstlisting}
# estrae le variabili di risposta e le unisce ai livelli dei fattori
python3 scripts/measurements/analyze_doe.py \
        --data-table 2-DoE/data_table.csv --out 2-DoE/results.csv

# statistiche descrittive e confronti fra configurazioni
python3 scripts/measurements/report_doe.py \
        --results 2-DoE/results.csv --out 2-DoE/REPORT.md

# figure
python3 scripts/measurements/plot_doe.py
\end{lstlisting}
